\section{Discussion}

This study provides a system-level evaluation of Liquid/CfC-style neural networks for cancer multi-omics prediction. The analysis combines a deep BRCA external-validation study with a broader five-task cancer-internal benchmark. The main conclusion is nuanced. Liquid/CfC did not universally dominate all baselines, especially after the baseline suite was expanded. After increasing the BRCA training data by adding METABRIC, the multimodal Liquid/CfC model was highly competitive on CPTAC 2020, but elastic-net logistic regression and MLP mRNA-marker baselines slightly exceeded it in macro F1.

The results highlight five points. First, breast cancer molecular subtype prediction is strongly driven by subtype-specific expression programs. This explains why mRNA marker models remained highly competitive across external cohorts. Second, multi-omics fusion is not automatically superior to expression-only modeling. It appears most useful when the external cohort has genuinely aligned multi-omics modalities, but even then it must be compared with strong expression baselines. Third, reducing Liquid/CfC capacity through small-Liquid/CfC lowered the parameter count substantially but did not reliably improve external generalization. Fourth, increasing mRNA training scale by adding SCAN-B and SMC 2018 did not improve the best CPTAC external performance, suggesting that cohort and platform shifts can outweigh the benefit of larger sample size. Fifth, the multi-cancer benchmark showed that endpoint difficulty and biological context strongly influence which model family performs best.

These findings argue against overstating Liquid/CfC as a universally superior classifier. A more defensible interpretation is that Liquid/CfC provides an input-adaptive cross-omics state-update mechanism. It is especially relevant when the scientific question concerns how different molecular layers update a latent disease-state representation, rather than simply maximizing subtype classification accuracy.

From a methodological perspective, the study also illustrates why external validation, extensive baseline testing, sample-alignment auditing, cross-validation, and task diversity matter for multi-omics neural networks. A model that appears promising under one internal split may not retain the same ranking across cohorts with different platforms, class distributions, and available modalities. Likewise, a model that is competitive for one endpoint may not be optimal for another. The core-aligned sensitivity analysis showed that removing sparse RPPA features did not turn Liquid/CfC into a universal winner, although it improved the Liquid/CfC ranking for COADREAD molecular subtype in the three-seed split analysis. Ten-fold CV further showed that the top model and runner-up were not significantly different after multiple-testing correction in any internal task. The strongest claim supported by the present results is not that Liquid/CfC replaces simpler baselines, but that it is a competitive and interpretable candidate for cross-omics state-update fusion in selected settings.

\subsection{Limitations}

Several limitations remain. CPTAC 2020 provided multimodal external validation but had only 122 labeled samples, so the strong multimodal Liquid/CfC result should be interpreted cautiously. SMC 2018 and SCAN-B provided independent expression validation but could not test complete mRNA+CNA+mutation fusion. The modality sequence used by Liquid/CfC is biologically motivated but not a true temporal sequence. The TCGA cancer-internal benchmark is sample-ID aligned, but the RPPA block is sparse; therefore, the full RPPA-containing benchmark should be interpreted together with the core-aligned sensitivity analysis. The ten-fold CV and noise analyses suggest that several model differences are small enough to be consistent with random fold-to-fold or perturbation-level variation. Finally, subtype labels can differ across studies due to platform, preprocessing, and annotation differences.

Another limitation is that the final BRCA marker panel was intentionally subtype-focused. This improves biological relevance for breast cancer molecular subtype prediction but limits conclusions about generic cancer-driver panels. The multi-cancer extension partly addresses this limitation, but it remains an internal-split benchmark rather than a multi-cohort external validation for each cancer type. Future studies should add external validation cohorts for UCEC, COADREAD, HNSC, KIRC, and PRAD, and should evaluate endpoints where genomic, epigenomic, proteomic, and transcriptomic layers provide complementary information, such as treatment response, recurrence risk, or survival.

\subsection{Conclusion}

Adding METABRIC to the training pool improved the viability of Liquid/CfC for external breast cancer subtype prediction. The original Liquid/CfC model outperformed the reduced small-Liquid model in most BRCA external settings and was highly competitive on CPTAC 2020 multimodal validation. Nevertheless, strengthened expression-marker baselines slightly exceeded Liquid/CfC on CPTAC macro F1, and larger mRNA training scale did not improve the best external performance. In the multi-cancer internal benchmark, the best model varied by endpoint, and ten-fold CV did not show statistically significant top-versus-runner-up differences after correction. The most appropriate conclusion is that Liquid/CfC is a promising but context-dependent cross-omics fusion architecture, not a universally dominant model.
